\section{Resolution map for the repository difficulties}

\subsection{What the new framework resolves}

The contribution is architectural rather than a claim that every difficult theorem in the
repository is now proved.  The following table separates resolved representation problems from
remaining mathematical producer problems.

\begin{longtable}{@{}p{0.27\textwidth}p{0.31\textwidth}p{0.34\textwidth}@{}}
\toprule
Difficulty & Resolution supplied here & Remaining boundary \\
\midrule
\endfirsthead
\toprule
Difficulty & Resolution supplied here & Remaining boundary \\
\midrule
\endhead
Quotients without decidable normal forms & Proof-relevant ideal-combination equality; operations transform witnesses. & Finding a witness may require a separate Gröbner or elimination theorem. \\
Localization with zero divisors & Equality carries a killing exponent; transitivity and congruence have explicit formulas. & Invertibility or annihilation witnesses must be supplied or produced. \\
Finite local-to-global gluing & Power-cover constructor plus explicit module gluing and uniqueness. & Arbitrary infinite covers are not internalized; use finite refinements or finite observations. \\
Scheme representation & Certified finite affine atlas with localization isomorphisms and cocycles. & No single completed internal universe of all schemes or prime ideals is claimed. \\
Sheaf condition & Exact equalizer for modules on finite principal covers. & General sheaf categories and derived functors require a larger coded library. \\
Projective geometry & Standard $\Proj$ charts and transitions are primitive-recursively generated. & Global projective existence theorems may still require nontrivial algebraic certificates. \\
Projective cohomology & Explicit finite \v{C}ech normal forms and contracting homotopy for $\cO(d)$. & Formal comparison with derived-functor cohomology is a separate extension. \\
Chart/branch nonuniqueness & Combine certified atlases with the prior finite-use relational framework; require invariant outputs. & Discontinuous exact selectors remain unavailable without stronger input. \\
Object-theory boundary & Exact list of \code{PR}/\code{PTm}/\code{Prf PRAAx} targets and quantifier-free invariants. & The final Lean implementation and compilation have not been completed in this archive. \\
\bottomrule
\end{longtable}

\subsection{Integration with Diophantine geometry}

The existing Diophantine packet represents rational projective points, heights, polynomial maps,
conics, local congruence obstructions, and finite descent.  The atlas layer extends that work in four
ways.

\begin{enumerate}[label=\arabic*.]
\item \textbf{Local equations.}  A projective variety can be dehomogenized on finitely many charts,
      and local algebraic manipulations can be glued or compared by localization witnesses.
\item \textbf{Rational maps with domains.}  A denominator-nonvanishing condition is represented by a
      principal open.  Several formulas for the same rational map are compared on overlaps rather
      than forced into one global syntactic expression.
\item \textbf{Line bundles and sections.}  Height constructions often use homogeneous forms or
      divisors.  Transition factors and local sections can now be represented and checked on a
      finite atlas.
\item \textbf{Local-to-global obstructions.}  Congruence, real, or $p$-adic local certificates can be
      attached to chart data, while exact algebraic descent witnesses handle the finite overlaps.
\end{enumerate}

A bounded rational-point search on a projective scheme proceeds as before: enumerate primitive
coordinate tuples up to a supplied height, check homogeneous equations, and quotient by projective
scaling.  The atlas does not create the height bound.  It provides the geometric operations needed
to state and verify local computations around the enumerated points.

\subsection{Faltings and Siegel}

The Faltings and Siegel packets correctly distinguish a finite cardinality statement from an
effective height cutoff.  The present framework sharpens the interface but does not remove that
obstruction.

A positive theorem can have the form:

\begin{quote}
Given a presented projective curve $X$, a height bound $B$, and a certificate proving that every
rational point has height at most $B$, the bounded enumeration returns a complete list; every local
chart computation and overlap equality is certified by the atlas.
\end{quote}

The bounded enumeration, chart conversion, and completeness-from-bound implication are primitive
recursive.  The production of $B$ for general high-genus curves is the missing arithmetic theorem.
A count bound, a finite-dimensional span bound, or a point-count estimate over finite fields is not a
substitute.

For Runge-type integral-point problems, explicit pole data and a supplied separation/height estimate
can be attached to the finite projective compactification.  This may simplify the exact bookkeeping
of points at infinity, divisors, and local functions, but the quantitative Runge inequality remains
the substantive producer.

\subsection{Mordell--Weil descent}

The Mordell--Weil packet has an explicit finite-height descent conditional on actual representatives
of the quotient by doubling.  A finite group scheme, isogeny, or Selmer computation can now be
presented chartwise with exact localization certificates.  This is useful for verifying the descent
maps and local images.

It still does not follow that a finite abstract Selmer class has a rational point representative.
The atlas framework can verify a supplied representative and its compatibility; it cannot create a
rational point from a purely cohomological class without an additional theorem.

\subsection{Puiseux expansions and singularities}

For exact polynomial coefficients, Newton polygons, finite truncations, and branch equations can be
encoded as finite algebraic certificates.  A singularity may be covered by finitely many candidate
charts or blow-up charts; relational finite-use composition can avoid an arbitrary global branch
label.

For arbitrary Cauchy coefficient names, the exact vanishing tests that distinguish singular strata
remain discontinuous.  The correct interfaces are:

\begin{itemize}
\item a supplied exact algebraic stratum or Newton-polygon certificate;
\item a derivative/separation witness giving a regular local branch;
\item finite root clusters and branch-invariant observations;
\item approximate solution germs with residual bounds.
\end{itemize}

The finite atlas does not turn approximation into exact equality.

\subsection{Analytic number theory packets}

Chebotarev, Siegel--Walfisz, Bombieri--Vinogradov, and related packets require uniform analytic
estimates, exceptional-zero control, field arithmetic, and explicit error terms.  Finite algebraic
geometry can organize number fields, finite covers, discriminants, and Frobenius data, but it does
not supply the missing estimates.  The no-universal-evaluator and finite-observation disciplines
continue to apply.

\subsection{Further algebraic-geometry targets}

The framework gives a clear route to several additional iPRA projects.

\begin{enumerate}[label=\arabic*.]
\item \textbf{Blow-ups of finitely presented affine schemes.}  Use the standard finite charts of the
      Rees algebra and localization transition maps.
\item \textbf{Grassmannians.}  Use Pl\"ucker coordinates, finitely many standard minors, and explicit
      transition matrices.
\item \textbf{Finite morphisms and normalization certificates.}  Present the target algebra as a
      finite module with multiplication tables and verify integral equations.
\item \textbf{B\'ezout in zero-dimensional certified intersections.}  Combine resultants or a
      homogeneous Gröbner certificate with finite quotient-rank calculations.
\item \textbf{Riemann--Roch for explicitly presented curves.}  First formalize divisor/line-bundle
      \v{C}ech complexes and finite-dimensional normal forms; any general theorem still needs
      quantitative algebraic inputs.
\item \textbf{Cohomology of toric varieties.}  Fans are finite combinatorial objects, affine toric
      charts are monoid algebras, and line-bundle cohomology often reduces to finite lattice-region
      calculations.
\item \textbf{Finite descent for vector bundles.}  Transition matrices and principal-cover gluing
      already provide the base layer for explicit projective-module calculations.
\end{enumerate}

The common theme is that the geometric theorem should expose a finite atlas, a finite complex, or a
finite certificate rather than conceal infinitary semantics behind an unexamined object code.
